\section{Introduction}

A grower deciding when to send pickers into a greenhouse needs to know, for
each row, how many fruit are ready now and how many will be ready next week.
Harvest scheduling and short-term yield forecasting both rely on these
per-stage counts. In the longer term, an automated harvesting robot can use the
same counts to decide which fruit to pick \cite{kootstra2021harvesting, yaguchi2016tomato}.
Cameras can provide these counts because tomato maturity is judged almost
entirely from colour, which can be measured from an image.

The deployment device affects both latency and reliability. Greenhouse rows
often have poor network coverage. Sending every frame to a cloud server also
adds a network round trip to each prediction, making latency higher and less
predictable than on-device inference \cite{kang2017neurosurgeon, shi2016edge}.
A smartphone provides a camera and a dedicated neural accelerator. This work
therefore assumes that inference runs on such an edge device and uses the camera
frame rate as the time budget for the detector. Running convolutional detectors
at camera frame rate on a smartphone has only recently become practical. Apple
introduced its on-device inference framework, Core ML, in June 2017, and its
first smartphone Neural Engine, in the A11 chip, in September of the same year
\cite{apple2017ios11, apple2017iphonex}.

Public datasets with annotated greenhouse tomatoes and maturity labels have
existed since 2020 \cite{laborotomato2020, tsironis2020tomatod}. Detectors of
the YOLO family \cite{redmon2016yolo, khanam2024yolov11} train on them in under
an hour, and studies published between 2020 and 2025 report mean average
precision (mAP) values that improve on one another by a few points
\cite{liu2020yolotomato, su2022seyolov3, li2023yolov5tomato, li2023mhsayolov8, camacho2023tomato}. In almost all of these studies, a
detector is trained on one dataset and its mAP is reported on that dataset's
validation split.

A counting error reported in this way depends on at least six choices. The
model is only one of them. First, annotation teams may use different meanings
of \emph{ripe}, so a model learns the colour boundary used in its training
dataset. Second, leaves and neighbouring fruit conceal parts of many fruit. The
construction of an occlusion test can therefore change its result. Third, a
count includes every detection above a confidence threshold. The number of
detections retained depends on the calibration of the model's confidence
scores, as well as its ability to rank fruit \cite{guo2017calibration}. Fourth,
lighting changes the colour used to determine the stage. Fifth, larger models
require more memory and inference time, which may exceed the resources of an
edge device. Sixth, an exported model can use different numerical precision,
preprocessing and hardware from the training checkpoint used for evaluation.

Each choice can change the reported result while the trained weights remain
fixed, yet papers usually report a single figure. Two papers can therefore
report the same counting error for the same crop while using different
measurement conditions. A speed measured on a development machine may also
differ on the intended deployment device.

This work asks: \emph{how much does the evaluation protocol contribute to a
reported tomato counting error, how much does the model contribute, and what
must be reported for the figure to be reproducible?} To answer these questions,
we vary the annotation convention, the occluder, the confidence threshold, the
lighting, the model size and the inference hardware one at a time. We keep the
detector fixed wherever possible and measure the effect of each change.

Previous work has studied individual parts of this problem, but has not
examined them together. Tomato maturity studies train and test on a single
dataset and report mAP
\cite{su2022seyolov3, li2023yolov5tomato, li2023mhsayolov8, camacho2023tomato}.
Detectors for tomatoes have been run on Android phones
\cite{zeng2023mobile, jia2025tseyolo}, on an NVIDIA Jetson board
\cite{liu2025gaeyolo}, and, for fruit detection without maturity stages, on the
Neural Engine of an iPhone \cite{kafi2025drone}; each of these reports the
speed of one or a few models on one platform. To our knowledge, this is the
first study that compares the two public tomato maturity datasets under one
label set, and the first to report maturity detection accuracy, counting error
and Neural Engine latency together for a complete YOLO11 model-size range (n
to x) on a smartphone.

This work makes six contributions:
\begin{enumerate}
\item It measures how well detectors transfer between Laboro Tomato and
tomatOD after both are mapped onto one three-stage label set, and explains the
asymmetric errors by measuring where each dataset places its colour boundaries.
\item It separates detection, stage assignment and counting under controlled
occlusion, and shows how the matching threshold and the occluder's appearance
change the result.
\item It shows that a fixed confidence threshold can reverse a cross-dataset
counting comparison in which mAP shows no difference.
\item It follows one trained model from FP32 PyTorch, through FP16 Core ML, to
a physical iPhone, and checks that the laptop and the smartphone produce the same
detections.
\item It trains all five YOLO11 sizes under one protocol and measures accuracy,
counting error, package size and latency on a laptop and on the smartphone. Accuracy
stops improving at the medium model (YOLO11m), while the largest model
(YOLO11x) only just meets the real-time budget, so the two sizes above YOLO11m
increase latency without improving accuracy.
\item It measures how much a change in exposure or white balance of the test
images reduces accuracy, and whether a wider colour augmentation during
training recovers the loss, with every configuration trained at three seeds so
that real effects can be separated from training noise
\cite{bouthillier2021variance}.
\end{enumerate}

The rest of the paper is organised as follows.
Section~\ref{sec:background} reviews tomato maturity detection, the public
datasets, how detectors are evaluated, and previous work on robustness,
occlusion and on-device inference, and identifies the gap this work addresses.
Section~\ref{sec:methods} describes the datasets, experimental protocols and
statistical analysis. Section~\ref{sec:results} reports the six experiments.
Sections~\ref{sec:discussion} to \ref{sec:conclusion} discuss the results,
describe the software artefact and its limitations, and conclude.
